\section{引言}

全景语义分割作为场景理解的核心任务，在自动驾驶、智能监控等领域具有广泛应用前景\cite{opencv_fisheye}。与传统透视图像相比，360°全景图像能够提供完整的周围环境信息，但其独特的拓扑结构和投影畸变也带来了新的技术挑战。

近年来，基于视觉大模型的分割方法取得显著进展，SAM2\cite{ravi2024sam2} 通过时序记忆机制实现了高质量的视频目标分割。OmniSAM\cite{omnisam2024} 将此思想扩展至全景域适应场景，将360°全景图划分为重叠的透视窗口序列，利用记忆库(memory bank)累积跨窗口的上下文信息。然而，该方法存在两个核心问题：\textbf{(1) 训练效率}：大量历史记忆Token参与交叉注意力计算，导致训练时间和显存占用显著增加；\textbf{(2) 边界不连续}：尽管全景图在物理上首尾相连(0°$\equiv$360°)，但模型将x=0和x=W视为断裂边界，导致跨边界物体(如交通标志)的分割质量下降。

针对上述问题，本文提出一种高效的全景分割框架，包含三项核心创新：

\begin{itemize}
\item \textbf{记忆Token压缩(MTC)}：通过卷积聚合空间相邻的冗余记忆，在保持语义完整性的前提下将Token数量压缩8倍(32768$\rightarrow$4096)。
\item \textbf{稀疏记忆检索(SMR)}：基于查询-键相似度动态选择重要记忆子集，避免无关历史信息干扰当前预测。
\item \textbf{无缝循环记忆(SCM)}：在推理阶段对全景图进行循环填充，使滑动窗口能够感知跨边界的连续上下文，修复边界区域的分割断裂。
\end{itemize}

我们在Cityscapes$\rightarrow$DensePASS无监督域适应基准上进行了系统实验验证。\textbf{效率方面}，MTC+SMR组合在5个随机种子上平均达到60.00$\pm$0.05\% mIoU，相比基线(60.22$\pm$0.02\%)仅损失0.24\%，但注意力计算时间降低55.8\%，GPU显存节省1.2GB。\textbf{精度方面}，SCM在推理阶段提升整体mIoU 0.33\% (60.51$\rightarrow$60.84)，其中交通标志类别从26.07\%提升至31.11\%(+5.04\%)，验证了边界连续性建模对边界敏感类别的显著改善。多随机种子统计验证表明我们的方法具有良好稳定性和可重复性。

本文的主要贡献总结如下：

\begin{enumerate}
\item 提出记忆Token压缩与稀疏检索的协同优化方案，实现8倍压缩率和55.8\%时间降低，为大规模全景分割提供高效训练范式。
\item 设计无缝循环记忆机制，首次系统性地解决全景分割中的边界不连续问题，对交通标志等边界敏感类别提升显著。
\item 通过5-seed多随机种子实验建立统计显著性验证，为全景域适应研究提供可靠的性能基准。
\end{enumerate}
